\documentclass[../../../include/open-logic-section]{subfiles}

\begin{document}

\olfileid{sth}{z}{pairs}
\olsection{配对}

接下来要考虑的公理如下：

\begin{axiom}[配对公理]
对任意集合 $a, b$，集合 $\{a, b\}$ 存在。\[
	\forall a \forall b \exists P \forall x (x \in P \liff (x = a \lor x = b))
\]
\end{axiom}

下面说明如何用迭代概念来证成这条公理。假设 $a$ 在阶段 $S$ 可用，$b$ 在阶段 $T$ 可用。令 $M$ 为阶段 $S$ 和 $T$ 中较晚的那一个。那么由于 $a$ 和 $b$ 在阶段 $M$ 都可用，集合 $\{a,b\}$ 就是一个可能的汇集，在 $M$ 之后的任何阶段（即更大的那一个阶段之后）都是可用的。

但是且慢！为何要假定 $M$ \emph{之后}还有阶段呢？如果没有，那么我们的证成就会失败。所以，为了证成配对公理，我们必须在 \olref[sth][z][story]{sec} 中所讲的故事里补充另一条原则，即：

\begin{enumerate}
	\item[] \stagessucc. 没有最后一个阶段。
\end{enumerate}

这条原则本身合理吗？Shoenfield 的故事中并没有\emph{明确}说明没有最后一个阶段。不过，即使它（严格来说）是我们故事的一个额外补充，它与“集合是分阶段形成的”这一基本思想也是吻合的。在接下来的讨论中，我们将直接接受它。因此，我们也将接受配对公理。

凭借这条新公理，我们可以证明更多集合的存在性。例如：

\begin{prop}\ollabel{prop:pairsconsequences}
对任意集合 $a$ 和 $b$，下列集合存在：
	\begin{enumerate}
		\item\ollabel{singleton} $\{a\}$
		\item\ollabel{binunion} $a \cup b$
		\item\ollabel{tuples} $\tuple{a, b}$
	\end{enumerate}
\end{prop}

\begin{proof}
\olref{singleton}. 由配对公理，$\{a, a\}$ 存在，根据外延性公理，它就是 $\{a\}$。

\olref{binunion}. 由配对公理，$\{a, b\}$ 存在。进而由并集公理，$a \cup b = \bigcup
\{a, b\}$ 存在。

\olref{tuples}. 由 \olref{singleton}，$\{a\}$ 存在。由配对公理，$\{a,
b\}$ 存在。再由配对公理，$\{\{a\}, \{a, b\}\} = \tuple{a, b}$ 存在。
\end{proof}

\begin{prob}
证明：对任意集合 $a, b, c$，集合 $\{a, b, c\}$ 存在。
\end{prob}

\begin{prob}
证明：对任意集合 $a_1, \ldots, a_n$，集合 $\{a_1, \ldots,
a_n\}$ 存在。
\end{prob}

%\begin{proof} By two applications of Pairs, $\{\{a_1, a_2\}, \{a_1,
%   a_3\}\}$ exists. By Union and Extensionality, $\bigcup \{\{a_1,
%   a_2\}, \{a_1, a_3\}\} = \{a_1, a_2, a_3\}$ exists. Repeat this
%   trick as often as necessary. \end{proof}

\end{document}